\section{The FTC as an explicit finite-output theorem}
\label{sec:ftc-finite-output}

The preceding theorem is mathematical, but its statement must still specify
what is constructed and what proves termination.  We give that specification
here.  This also states the theorem with an independently supplied derivative
$g$: its conclusion is $\int_a^b g=f(b)-f(a)$, not merely an assertion
that some unspecified derivative exists.

\subsection{Orders, derivatives, and hypotheses}
For raw computations $X_n=[x_n^-,x_n^+]$ and $Y_n=[y_n^-,y_n^+]$, use
\[
 X\le Y\quad\Longleftrightarrow\quad
 x_n^-\le y_m^+\quad\text{for every }n,m\in\N.
\]
This is a predicate on rational endpoints.  It is compatible with arithmetic
and equivalence: two computations are equivalent exactly when both orders
hold.  Transitivity follows by choosing a sufficiently narrow intermediate
box if a positive rational gap were to separate the other two.  No point
in an infinite intersection is selected.  Bounds below use this order.

For a function computation $f$, the notation $f'(x)=g(x)$ means the
rational-input derivative relation
\[
 \forall\varepsilon\in\Q_{>0}\ \exists\delta\in\Q_{>0}\ 
 \forall h\in\Q,\qquad
 0<|h|<\delta\ \Longrightarrow\
 |f(x+h)-f(x)-h g(x)|\le\varepsilon |h|,
\]
with increments restricted to the declared domain.  This relation does not
choose a derivative by an existence principle.  The secant algorithm
constructs the candidate; its shrinking-gap proof proves validity, and
the earlier secant sandwich proves this relation.  Uniqueness is proved
by using any sufficiently small nonzero rational increment in the two
relations and then taking arbitrarily small rational tolerances.

Fix rational $A<a<b<B$.  Suppose $p,q$ are specified number-valued
computations on the rational points of $[A,B]$, both convex in the finite
rational sense.  For every rational $x\in[a,b]$, assume their centered
secant gaps shrink to zero.  Explicitly, for $v=p$ and $v=q$, put
\[
 h_0(x)=\tfrac12\min(x-A,B-x),\qquad h_n(x)=2^{-n}h_0(x).
\]
The required shrinking statement is
\[
 \forall x\in[a,b]\cap\Q\ \forall\varepsilon\in\Q_{>0}\
 \exists n,\qquad
 \frac{v(x+h_n)-2v(x)+v(x-h_n)}{h_n}<\varepsilon.
\]
It proves termination of the rational-box precision search already
specified in Definition~\ref{def:computed-derivative}.  The algorithm
searches for a rational box width, not for an undecidable comparison of
computed values.  The shrinking proof is not consulted at run time.
No modulus uniform in $x$ is assumed.

Let $f$ be a specified computation with $f(x)\sim p(x)-q(x)$ on this
domain.  Then $d(x)=Dp(x)-Dq(x)$ is a valid computed derivative of $f$.
If $g$ is any other specified computation satisfying $f'(x)=g(x)$ at
every rational $x\in[a,b]$, uniqueness gives $g(x)\sim d(x)$.  The
following output theorem therefore applies to $g$ as well as to $d$.

\begin{theorem}[Finite-output fundamental theorem]
\label{thm:ftc-finite-output}
Under the hypotheses just given, the independent cell-range integral
construction produces a valid raw computation $J$ of $\int_a^b g$.
For every stage of $J$ and every stage of the independently computed
endpoint difference $f(b)-f(a)$, their rational intervals overlap.
Consequently
\[
 \boxed{\int_a^b g(x)\,dx\sim f(b)-f(a).}
\]
The construction below gives a terminating precision schedule and a
finite bound for every tagged rational sum.
\end{theorem}

\subsection{Rational cell ranges and the error schedule}
Put $L=b-a$.  Secant ordering proves that $Dp$ and $Dq$ are nondecreasing.
Coarse enclosures of their four endpoint values give a rational $V\ge0$
satisfying
\[
 V\ge Dp(b)-Dp(a)+Dq(b)-Dq(a).
\]
For example, add the upper bounds at $b$ minus the lower bounds at $a$
for the two computations.  This finite quantity is already nonnegative.
Alternatively outer secants provide such a bound before evaluating the
endpoint derivatives.

Choose a rational partition $P$ and compute, at every endpoint $x_i$,
\[
 Dp(x_i)\in[\alpha_i^p,\beta_i^p],\qquad
 Dq(x_i)\in[\alpha_i^q,\beta_i^q]
\]
to widths at most $\theta>0$.  The actual rational ranges used for $g$
on the cell $[x_i,x_{i+1}]$ are
\[
 l_i=\alpha_i^p-\beta_{i+1}^q,\qquad
 u_i=\beta_{i+1}^p-\alpha_i^q.
\]
Monotonicity and $g\sim Dp-Dq$ prove $l_i\le g(t)\le u_i$ for all
rational $t$ in that cell.  Set
\[
 J(P,\theta)=\left[\sum_i l_i\Delta x_i,\sum_i u_i\Delta x_i\right].
\]
Every endpoint of this interval is rational.  Finite telescoping gives
\[
 \boxed{\operatorname{width}J(P,\theta)
           \le\mesh(P)V+4\theta L.}
\]
The four error terms per cell are precisely the four derivative endpoint
enclosures used above.  For one convex primitive alone, the constant is
$2\theta L$ instead.

For requested width $\varepsilon>0$, take an equal partition with integer
$N$ satisfying $LV/N\le\varepsilon/4$, and take
$\theta\le\varepsilon/(16L)$.  An integer meeting this condition is
found by rational arithmetic; for $V=0$ any positive $N$ suffices.
There are only $2(N+1)$ terminating derivative precision searches at this
stage.  The resulting interval has width at most $\varepsilon/2$.
Use $\varepsilon=2^{-n}$ successively and intersect the finitely many
output intervals obtained so far.

Those intersections are nonempty independently of any primitive endpoint:
for two partitions, take their finite common refinement.  On each refined
cell choose a rational sample.  Its computed value is bounded by both
parent cell ranges; adding the inequalities shows that each lower sum is
at most the other upper sum.  Thus the two rational intervals overlap.
A finite family of pairwise overlapping rational intervals has nonempty
intersection, by finite maxima and minima.  This proves integral validity
without defining the integral through $f(b)-f(a)$.

\subsection{The endpoint identity is a separate proof}
For each cell, the convex secant inequalities give
\[
 Dp(x_i)\Delta x_i\le p(x_{i+1})-p(x_i)
                         \le Dp(x_{i+1})\Delta x_i,
\]
and similarly for $q$.  Subtract the second pair from the first, then
apply the rational endpoint enclosures.  The result is
\[
 l_i\Delta x_i\le f(x_{i+1})-f(x_i)\le u_i\Delta x_i.
\]
Adding telescopes to $f(b)-f(a)$.  The raw-order definition shows that
every enclosure of that difference overlaps $J(P,\theta)$, and hence
every finite-prefix intersection of these intervals.  This proves the
stated equivalence after, and separately from, integral validity.

For any rational tag $\xi_i\in[x_i,x_{i+1}]$, monotone derivative ranges
also give
\[
 \left|f(b)-f(a)-\sum_i g(\xi_i)\Delta x_i\right|
                 \le\mesh(P)V.
\]
Replacing the tags' values by rational numbers $r_i$ with proved errors
$|r_i-g(\xi_i)|\le\eta$ yields the fully quantitative formula
\[
 \boxed{\left|f(b)-f(a)-\sum_i r_i\Delta x_i\right|
                     \le\mesh(P)V+\eta L.}
\]
This permits a different evaluator for $g$, unrelated to the secant
evaluator, once pointwise derivative uniqueness has identified it.

\subsection{The reverse direction and the boundary of the theorem}
If a supplied nondecreasing computation $g$ is continuous at each rational
point of an interval, its monotone rectangle integral constructs
$F(x)=\int_a^x g$.  Ordered cell averages prove convexity of $F$.
At an interior rational $x$, its left and right secants lie between
$g(x-h)$ and $g(x+h)$.  Pointwise continuity makes the gap shrink, so
the secant derivative computation terminates and $DF(x)\sim g(x)$.
Thus accumulation constructs a primitive as well as recovering an
already supplied one.  One-sided endpoint versions require their
explicit domain and endpoint conditions.

The convex-difference hypothesis is a sufficient computational construction,
not a claim that every useful function must be presented that way.
Subsequent series and local integral constructions establish their own
finite increment bounds and transfer their identities through proved
errors.  A formal proof must retain these bounds, rather than infer an
FTC from pointwise rational differentiation alone.  Explicit jumps are
handled separately by their delta contributions; they are not silently
admitted by the present theorem.
